%% 01_introduction.tex
%% Introduction section

\section{Introduction}
\label{sec:introduction}

Academic writing requires careful attention to formatting and structure.
This template provides a foundation for preparing papers for ICSE submissions
using the ACM sigconf format, incorporating best practices from the academic community.

\subsection{Background}

The International Conference on Software Engineering (ICSE) is a premier venue
for software engineering research. A well-designed template helps authors focus
on content while ensuring consistent formatting that meets ACM publication standards.

\subsection{Features}

This template includes several useful features:

\begin{itemize}
    \item Two-column layout based on ACM sigconf format
    \item Mathematical typesetting with theorem environments
    \item Code listing support with syntax highlighting
    \item Algorithm description capabilities
    \item Automatic bibliography management with ACM Reference Format
\end{itemize}

\subsection{Mathematical Examples}

The template supports various mathematical constructs.
For example, a simple equation:

\begin{equation}
    E = mc^2
    \label{eq:einstein}
\end{equation}

And a more complex aligned equation:

\begin{align}
    \nabla \cdot \mathbf{E} &= \frac{\rho}{\varepsilon_0} \\
    \nabla \cdot \mathbf{B} &= 0 \\
    \nabla \times \mathbf{E} &= -\frac{\partial \mathbf{B}}{\partial t} \\
    \nabla \times \mathbf{B} &= \mu_0 \mathbf{J} + \mu_0 \varepsilon_0 \frac{\partial \mathbf{E}}{\partial t}
\end{align}

\begin{theorem}
\label{thm:example}
For any real numbers $a$ and $b$, $(a + b)^2 = a^2 + 2ab + b^2$.
\end{theorem}

\begin{proof}
By direct expansion of $(a + b)(a + b)$.
\end{proof}

\subsection{Code Listing Example}

The template supports code listings with syntax highlighting:

\begin{lstlisting}[language=Python, caption={Example Python code}]
def hello_world():
    """Print a greeting message."""
    print("Hello, World!")
    return True

if __name__ == "__main__":
    hello_world()
\end{lstlisting}

\subsection{Algorithm Example}

Algorithms can be described using the algorithm2e environment:

\begin{algorithm}[t]
\caption{Example Algorithm}
\label{alg:example}
\KwIn{Input data $X$}
\KwOut{Processed result $Y$}
Initialize $Y \leftarrow \emptyset$\;
\For{each element $x$ in $X$}{
    Process $x$\;
    Add result to $Y$\;
}
\Return{$Y$}
\end{algorithm}

\subsection{Paper Organization}

The remainder of this paper is organized as follows.
Section~\ref{sec:conclusion} concludes the paper with a summary
of our contributions and future work directions.
